\chapter{Related Work}

This chapter reviews the work that motivates the study design rather than
surveying all LLM evaluation research. It is organized around six themes:
algorithm auditing; LLM behavior, deployment, and access surfaces; safety and
refusal calibration; control and factuality tasks; the reliability of LLM
measurement; and the positioning of this thesis. Each theme closes by stating
what prior work leaves open for an access-surface audit.

The empirical study groups its prompts into six behavioral families that
Chapter 3 formalizes as Categories A--F. To keep this chapter self-contained, the
work below is described by behavioral function (over-refusal, refusal-expected
guardrailing, boundary safety framing, instruction following, factual control,
and socially sensitive inference) rather than by category letter.

\section{Algorithm Auditing}

Algorithm auditing studies opaque systems through controlled inputs and
observable outputs. Sandvig et al.\ \cite*{sandvig_auditing} describe it as a
method for detecting discrimination and differential treatment on internet
platforms when internal mechanisms are unavailable. The logic transfers directly
to proprietary LLM systems: the evaluator cannot see hidden instructions, tool
policies, routing, or moderation layers, but can still compare outputs collected
under documented conditions.

Empirical audits show why the deployed product, not an abstract model, is the
right unit of analysis. Buolamwini and Gebru's Gender Shades study
\cite*{gender_shades} measured commercial gender-classification products and
exposed large intersectional accuracy gaps that aggregate scores had hidden;
Raji and Buolamwini \cite*{raji_actionable} then showed that publishing such
product-level audits changed vendor behavior; and Obermeyer et al.\
\cite*{obermeyer_bias} found racial bias in a widely deployed healthcare risk
tool by comparing its outputs against outcome data. The lesson taken here is not
the domain but the unit: an audit should test the deployed system and define its
boundary explicitly, rather than evaluate a nominal model in isolation.

That boundary problem is central to accountability work. Raji et al.\
\cite*{raji_gap} propose an end-to-end framework in which evidence, system scope,
and deployment assumptions are documented together. Selbst et al.\
\cite*{selbst_abstraction} warn that technical evaluations fail when they
abstract away the social and institutional context of deployment, and Suresh and
Guttag \cite*{suresh_guttag} show that harms can enter at many points of the
machine-learning life cycle, including deployment. A consumer app is therefore
not a transparent wrapper around an API endpoint; it is part of the deployed
condition a user encounters. What this literature does not supply is a method for
auditing \emph{generative, non-deterministic} text systems, where the same input
can yield different outputs; that measurement problem is taken up in
Section~2.5.

\section{LLM Behavior, Deployment, and Access Surfaces}

Modern LLM behavior is shaped by far more than pretraining. GPT-3 \cite*{gpt3}
established broad prompt-based task performance, and InstructGPT
\cite*{instructgpt} showed that instruction tuning with human feedback reshapes
what a model will and will not do. Anthropic's helpful, honest, and harmless
assistant work \cite*{helpful_harmless} and Constitutional AI
\cite*{constitutional_ai} make the same point for safety: assistant behavior is a
trained, system-level property, not raw next-token continuation. At a broader
level, Bommasani et al.\ \cite*{foundation_models} and Weidinger et al.\
\cite*{weidinger_risks} argue that the risks of foundation models depend on how
they are adapted and deployed. If behavior is a property of the deployed system,
then the deployed system is what an audit must measure.

Benchmarks are the standard instrument for that measurement. MMLU \cite*{mmlu},
GSM8K \cite*{gsm8k}, HumanEval \cite*{codex}, BIG-bench \cite*{bigbench}, and HELM
\cite*{helm} standardized prompt suites, metrics, and collection settings, and
HELM in particular treats scenario design and collection conditions as part of
evaluation transparency. But these suites are almost always run through
programmatic APIs. That makes them reproducible, yet it leaves a gap: whether the
behavior they measure is the behavior a user receives through a consumer
application is rarely tested.

A small but growing literature treats the access surface itself as an empirical
variable. The closest precedent is Dual Standards by Lipphardt et al.\
\cite*{dualstandards}, which provides empirical evidence, not merely a conceptual
argument, that moderation behavior differs between API and WebUI access for the
same providers. Notably they find the WebUI \emph{more} restrictive, the opposite
direction to the result reported in this thesis, which is precisely why the claim
here is only that the surfaces differ, not that they differ in a fixed direction.
Two further lines of work make the deployed app a moving, opaque target. Chen et
al.\ \cite*{chatgpt_drift} show that the behavior of the same named ChatGPT
service can change substantially within months, so the app is not even
temporally stable. And Zhang et al.\ \cite*{prompt_extraction} demonstrate that
consumer systems such as ChatGPT and Claude wrap the model in hidden system
prompts that can be partially extracted, which is direct evidence for the hidden
instruction layer this thesis treats as part of the deployment stack. What
remains missing is a direct, paired, evidence-backed comparison of a bare API
condition against the official consumer app on one fixed diagnostic prompt bank.
That is the gap this thesis addresses.

\section{Safety and Refusal Calibration}

Safety benchmarks supply the high-risk and boundary prompts the bank needs.
Do-Not-Answer \cite*{donotanswer} curates prompts a responsible model should
decline; HarmBench \cite*{harmbench} standardizes red-teaming and refusal
evaluation; AILuminate \cite*{ailuminate} organizes product-level risk and
reliability categories; and moderation classifiers such as Llama Guard
\cite*{llama_guard} show that a provider can place a safety filter between the
user and the model, which from a black-box view is part of the observed
deployment stack.

Refusal benchmarks sharpen how a non-answer should be coded. SORRY-Bench
\cite*{sorrybench}, StrongREJECT \cite*{strongreject}, and WildGuard
\cite*{wildguard} evaluate refusal, harmfulness, and jailbreak behavior from
different angles. The biological-research RefusalBench of Weidener et al.\
\cite*{refusalbench_bio} is the key cautionary reference: it shows that a single
refusal rate can misrank systems when safe, useful engagement is possible. (A
separate 2026 paper also uses the name ``RefusalBench''; the cautionary result
used here is specifically the biological-prompts study.) This thesis therefore
does not treat every non-answer alike; it distinguishes full refusal, partial
refusal, and safe redirection.

Over-refusal work motivates the benign half of the design. XSTest \cite*{xstest}
targets exaggerated safety on clearly safe prompts that contain superficially
risky words, and OR-Bench \cite*{orbench} scales over-refusal evaluation. These
sources support the over-refusal prompts in the bank. The bank also contains
refusal-expected prompts, with high-risk raw text redacted in shared artifacts,
and answerable safety-sensitive prompts where safe framing is expected but
blanket refusal is not. The aim is not to maximize attack success but to compare
refusal posture and safe redirection across access surfaces.

Red-teaming and jailbreak research is background rather than a prompt source.
Ganguli et al.\ \cite*{ganguli_redteam}, Perez et al.\ \cite*{redteam_lms}, Zou
et al.\ \cite*{universal_attacks}, and Chao et al.\ \cite*{pair} probe safety
adversarially. This thesis instead uses plain diagnostic prompts so that any
resulting differences can be read as deployment-surface effects rather than
jailbreak artifacts.

\section{Control and Factuality Tasks}

The bank also needs controls that are not primarily about safety. IFEval
\cite*{ifeval} motivates instruction-following items with explicit, checkable
constraints that can be scored without subjective judgment. It also exposes a
practical app-evaluation hazard: rendered markdown, visual separators, or
transcription artifacts can mimic format failures unless screenshots are checked
carefully, which is why this thesis codes format-compliance divergences
conservatively.

Factuality supplies a second control family, and the relevant work fills three
distinct roles. Short-answer factuality benchmarks, SimpleQA \cite*{simpleqa} and
TruthfulQA \cite*{truthfulqa}, define when a model should simply be correct.
Hallucination studies, including Maynez et al.\ \cite*{maynez_factuality}, the
survey of Ji et al.\ \cite*{hallucination_survey}, FActScore \cite*{factscore},
and SelfCheckGPT \cite*{selfcheckgpt}, motivate separating unsupported memory
answers from evidence-supported ones. Retrieval work, including dense passage
retrieval \cite*{dpr}, retrieval-augmented generation \cite*{rag}, WebGPT
\cite*{webgpt}, GopherCite \cite*{verified_quotes}, and retrieval-augmented
dialogue \cite*{blenderbot2}, shows that adding visible sources changes the
factual support a user receives. Because the consumer app can display sources
while the bare API cannot, visible app sources are coded here as a
deployment-surface signature rather than dismissed as cosmetic presentation.

Socially sensitive inference is the final control family, used as a negative
control. Work on bias in embeddings and generated text
\cite*{caliskan_bias,debiasing_embeddings,sheng_bias,realtoxicityprompts} shows
that models absorb demographic associations, and Blodgett et al.\
\cite*{blodgett_survey} argue that bias claims require explicit normative
grounding. Most directly, BBQ \cite*{bbq} is constructed so that the correct
response is often to decline an unsupported demographic inference. This thesis
therefore codes ``cannot determine from the information given'' as a correct
answer rather than a refusal, so that these controls do not inflate the refusal
result.

\section{Reliability of LLM Measurement}

Auditing a generative system raises measurement problems that a classification
audit does not. First, outputs are non-deterministic. Yuan et al.\
\cite*{llm_nondeterminism} show that the same prompt can produce different
responses across repeated runs and across hardware, from floating-point and
batching effects alone, before any deliberate sampling. An audit that collects a
single response per condition can therefore mistake run-to-run noise for an
access-surface difference, which is why this thesis repeats the API condition and
reports run-to-run stability. Second, scaling the coding of open-ended outputs
invites automated judging: MT-Bench \cite*{mtbench} and G-Eval \cite*{geval} use
strong models as evaluators, but their documented sensitivity to prompt wording
and answer position means automated labels need human checking rather than blind
trust. Third, when human coding is used, agreement must be quantified. Cohen's
\cite*{cohen_kappa} and Fleiss's \cite*{fleiss_kappa} kappa coefficients, with
McHugh's \cite*{mchugh_kappa} guidance on interpreting kappa magnitudes, are the
standard instruments. This thesis combines all three ideas, namely repetition for
non-determinism, automated coding checked against an independent blind second
coder, and kappa-based inter-rater reliability, so that a reported divergence is
unlikely to be an artifact of stochasticity or coder subjectivity.

\section{Positioning of This Thesis}

Taken together, the cited work establishes three points and leaves one open. It
establishes that LLM behavior is a property of the deployed system rather than
the bare model, that the access surface can change that behavior, and that
measuring a generative system requires controlling for non-determinism and coder
subjectivity. What it does not yet provide is a compact, evidence-backed, paired
audit that holds a fixed diagnostic prompt bank constant while directly comparing
a bare OpenAI API condition with the official ChatGPT app, and that codes only
observable differences with screenshot-backed evidence and reliability controls.
That is the contribution of this thesis.

The same framing also bounds the claims. The study can show that observable
outputs differ under documented access conditions, but it cannot identify the
hidden system prompt, retrieval policy, tool-routing decision, or moderation
component that produced a given app output. Naming those components is left to
future work; the central implication that survives is narrow and practical, that
LLM evaluations should name the access surface they measure and validate app
behavior directly rather than assuming it from API results.
